\documentclass[11pt]{amsart}
\usepackage{amsmath,amssymb,amsthm,booktabs}
\usepackage[margin=1.15in]{geometry}
\usepackage{hyperref}
\newtheorem{theorem}{Theorem}
\newtheorem{proposition}[theorem]{Proposition}
\newtheorem{lemma}[theorem]{Lemma}
\newtheorem{corollary}[theorem]{Corollary}
\theoremstyle{definition}
\newtheorem{remark}[theorem]{Remark}
\newcommand{\Prob}{\mathbb{P}}
\newcommand{\E}{\mathbb{E}}
\newcommand{\tmix}{\tau_1^{(2)}}
\newcommand{\hE}{\widehat{\E}}
\newcommand{\hw}{\widehat{w}}
\newcommand{\hpi}{\widehat{\pi}}
\newcommand{\hI}{\widehat{I}}
\newcommand{\hA}{\widehat{A}}
\newcommand{\hX}{\widehat{X}}
\newcommand{\hM}{\widehat{M}}
\newcommand{\hS}{\widehat{S}}
\newcommand{\hF}{\widehat{F}}
\newcommand{\lra}{\leftrightarrow}
\newcommand{\one}{\mathbf{1}}

\title[Lumping and the mixing time]{Lumping and the mixing time of reversible chains with $n$ states}
\author{Yitzchak Shmalo}
\address{Einstein Institute of Mathematics, The Hebrew University of Jerusalem, Jerusalem, Israel}
\date{November 2026}
\subjclass[2020]{60J10, 05C81}
\keywords{reversible Markov chain, lumping, mixing time, relaxation time, hitting time}

\begin{document}
\begin{abstract}
Aldous and Fill asked whether lumping states of a reversible Markov chain can increase the stopping-rule mixing time
$\tmix$ by more than an absolute constant factor. Hermon and Peres showed that for random walks on graphs of bounded
degree the factor can be of the order of the logarithm of the number of vertices, and no larger. For reversible
chains with $n\ge4$ states we show that the largest possible factor $K(n)$ satisfies
$\lfloor n/2\rfloor\le K(n)\le n-2$, and that $K(n)=\lfloor n/2\rfloor$ for $n\le6$ and for $n=8$. The upper bound
rests on two facts: lumping two states multiplies $\tmix$ by less than $2$, a constant that cannot be lowered, and the
relaxation time is at most $\tmix$. The lower bound comes from a random walk on a weighted comb. We also show that
lumping increases the mean hitting time of a set $A$ by a factor of at most $n-|A|$, and that this bound is sharp up to
a factor that tends to $2$ as $n\to\infty$.
\end{abstract}
\maketitle

\section{Introduction}\label{sec:intro}
Let $I$ be a finite set with $n\ge2$ elements and let $(w_{xy})_{x,y\in I}$ be a symmetric array of nonnegative
weights, loops $w_{xx}>0$ being allowed, such that the graph with edges $\{x,y\}$, $x\ne y$, $w_{xy}>0$, is connected.
Put $w_x=\sum_y w_{xy}$ and $w=\sum_x w_x$. The random walk on this network is the Markov chain $X$ on $I$ with
transition probabilities $p_{xy}=w_{xy}/w_x$; it is reversible with stationary distribution $\pi_x=w_x/w$, and every
irreducible reversible chain on $I$ is of this form. We write $T_y=\inf\{t\ge0:X_t=y\}$ for the hitting time of a
state $y$ and $T_A=\inf\{t\ge0:X_t\in A\}$ for the hitting time of a set $A$, so that $T_y=0$ for the walk started at
$y$, and $\E_\mu$ and $\Prob_\mu$ for expectation and probability for the chain started from the distribution $\mu$.

Given a map $f$ from $I$ onto a set $\hI$, the \emph{induced chain} $\hX$ is the random walk on $\hI$ with weights
\[
\hw_{\hat x\hat y}=\sum_{x\in f^{-1}(\hat x)}\ \sum_{y\in f^{-1}(\hat y)} w_{xy},\qquad \hat x,\hat y\in\hI .
\]
Its transition probabilities are $\Prob_\pi(f(X_1)=\hat y\mid f(X_0)=\hat x)$ and its stationary distribution is
$\hpi_{\hat x}=\pi(f^{-1}(\hat x))$; this is the chain of Aldous and Fill~\cite[Section~4.6.1]{AF}. In the language of
networks, the states with the same value of $f$ are glued together. We say that $\hX$ is obtained from $X$ by
lumping, and write $\hE$ and $\widehat{T}$ for its expectations and hitting times. If $g$ maps $\hI$ onto a third set,
the chain induced from $\hX$ by $g$ is the chain induced from $X$ by $g\circ f$, as both have the weights
$\sum w_{xy}$ over the preimages; so every lumping onto $q$ states of a chain with $m$ states is a composition of
$m-q$ lumpings of two states.

The parameter of~\cite[Chapter~4]{AF} that we study is
\[
\tmix=\max_{i\in I}\ \min_{U}\ \E_iU ,
\]
where the minimum is over all stopping times $U$, possibly randomized with a randomization independent of the future
of the chain, such that $X_U$ has distribution $\pi$ under $\Prob_i$. Lov\'asz and Winkler~\cite{LW95} showed that the
minimum is attained and found its value (see also~\cite[Theorem~9.4]{AF}), so that
\begin{equation}\label{eq:lw}
\tmix=\max_{i,j\in I}\bigl(\E_iT_j-\E_\pi T_j\bigr).
\end{equation}
For reversible chains $\tmix$ is equivalent, up to absolute constants, to the total variation mixing time of the
chain in continuous time~\cite{Aldous82}, \cite[Theorem~4.6]{AF}; see also~\cite{LW98}.

The maximal mean commute time $\tau^*=\max_{i,j}(\E_iT_j+\E_jT_i)$, the average hitting time
$\tau_0=\sum_j\pi_j\E_iT_j$ (which does not depend on $i$), the relaxation time $\tau_2=1/(1-\lambda_2)$, where
$\lambda_2$ is the second largest eigenvalue of $(p_{xy})$, and the Cheeger time constant $\tau_c$ do not increase
under lumping~\cite[Proposition~4.44]{AF}. For the mixing time the matter was left open:

\smallskip
\noindent\textbf{Open Problem 4.45 of~\cite{AF}.} \emph{Is the value of $\tmix$ in an induced chain bounded by $K$
times the value of $\tmix$ in the original chain, for some absolute constant $K$? For $K=1$?}
\smallskip

Hermon and Peres~\cite[Theorem~3(c)]{HP} answered both questions in the negative. They construct graphs $G_n$ of
bounded degree for which lumping together certain pairs of vertices increases the total variation mixing time of the
lazy simple random walk by a factor of order $\log|V(G_n)|$. Lumping does not increase the relaxation time, and the
mixing time of a lazy reversible chain is bounded below by a constant multiple of the relaxation time and above by a
constant multiple of the relaxation time multiplied by $\log(1/\min_x\pi_x)$; hence this is the largest possible
order in that setting, by the argument given in~\cite[Remark~1.8]{HP} for perturbations of the edge weights.

For a random walk on a weighted network $\min_x\pi_x$ can be arbitrarily small, and the bound in terms of
$\log(1/\min_x\pi_x)$ says nothing about the dependence on the number of states. This paper determines that dependence
up to a factor that tends to two as $n\to\infty$. Let
\[
K(n)=\sup\ \frac{\tmix(\hX)}{\tmix(X)} ,
\]
the supremum being over all reversible chains $X$ with at least two and at most $n$ states and all induced chains
$\hX$ of $X$. Put $U(2)=1$ and
\[
U(n)=\max_{2\le q\le n-1}\ \min\bigl(2^{\,n-q},\ q-1\bigr)\qquad(n\ge3).
\]

\begin{theorem}\label{thm:main}
For every $n\ge2$,
\[
\Bigl\lfloor\frac n2\Bigr\rfloor\ \le\ K(n)\ \le\ U(n)\ \le\ \max(1,n-2) .
\]
In particular $K(n)=\lfloor n/2\rfloor$ for $2\le n\le6$ and for $n=8$, and $3\le K(7)\le4$.
\end{theorem}

\begin{table}[tbp]
\centering
\begin{tabular}{lrrrrrrrrrrr}
\toprule
$n$ & 2 & 3 & 4 & 5 & 6 & 7 & 8 & 9 & 10 & 11 & 12\\
\midrule
$\lfloor n/2\rfloor$ & 1 & 1 & 2 & 2 & 3 & 3 & 4 & 4 & 5 & 5 & 6\\
$U(n)$ & 1 & 1 & 2 & 2 & 3 & 4 & 4 & 5 & 6 & 7 & 8\\
\bottomrule
\end{tabular}
\caption{The bounds of Theorem~\ref{thm:main} for small $n$.}\label{tab:bounds}
\end{table}

The upper bound combines two inequalities. The first concerns a single lumping of two states.

\begin{theorem}\label{thm:merge}
Let $\hX$ be obtained from $X$ by lumping two states and keeping the other states apart. Then
$\tmix(\hX)<2\,\tmix(X)$, and the constant $2$ cannot be replaced by a smaller one.
\end{theorem}

Applied $m-q$ times, Theorem~\ref{thm:merge} bounds the factor by $2^{m-q}$ for a lumping of a chain with $m$ states
onto $q$ states. The second inequality bounds the factor by $q-1$; it is a chain of four inequalities, three of which
are in~\cite{AF}. The fourth is the following comparison of the relaxation time with the stopping-rule mixing time, in
which the constant is one; whether this holds is asked in~\cite[Open Problem~4.48(b)]{AF}. Comparisons of the two
quantities up to absolute constants are in~\cite{Aldous82,LW98}.

\begin{lemma}\label{lem:relax}
Let $\lambda<1$ be a real eigenvalue of the transition matrix of an irreducible chain on a finite set, with a real
eigenfunction $g$, and let $i$ be a state at which $|g|$ is largest. If $U$ is a stopping time such that $X_U$ has the
stationary distribution $\pi$ under $\Prob_i$, then $\E_iU\ge 1/(1-\lambda)$. In particular $\tau_2\le\tmix$ for
every reversible chain.
\end{lemma}

Equality $\tau_2=\tmix$ holds for every chain with two states and for the simple random walk on a complete graph.

The lower bound in Theorem~\ref{thm:main} is approached along a family of random walks on a weighted comb with $2N+4$
states, described in Section~\ref{sec:comb}. Lumping each of $N$ teeth with a vertex of the spine, and the end of the
spine with a heavy leaf, multiplies $\tmix$ by a factor that tends to $N+2$. With $N=0$ only two states are lumped and
the factor tends to $2$, which shows that the constant in Theorem~\ref{thm:merge} is optimal. In these examples the
stationary distribution is very far from uniform, in agreement with the bound of~\cite{HP}; see
Remark~\ref{rem:pimin}.

The comb uses $N+1$ lumpings of two states and reaches the factor $N+2$. We do not know whether a composition of
several lumpings of two states can multiply $\tmix$ by more than one plus their number. If they cannot, the argument of
Section~\ref{sec:upper} bounds the factor for a lumping onto $q$ states by $\min(m-q+1,q-1)\le\lfloor m/2\rfloor$, and
then $K(n)=\lfloor n/2\rfloor$ for every $n$.

The mechanism behind the lower bound is a one-sided hitting time. Aldous and Fill observe~\cite[Section~4.6.1]{AF}
that, although commute times decrease under lumping, a mean hitting time may increase, and give an example on four
vertices where it doubles. The next proposition bounds the increase for hitting times of sets.

\begin{proposition}\label{prop:hit}
Let $\hX$ be the chain induced from $X$ by $f:I\to\hI$, let $\hA$ be a non-empty proper subset of $\hI$ and
$A=f^{-1}(\hA)$. Then for every $\hat x\in\hI\setminus\hA$
\[
\hE_{\hat x}\widehat{T}_{\hA}\ \le\ \sum_{y\in I\setminus A}\E_yT_A\ \le\ (n-|A|)\max_{y\in I}\E_yT_A .
\]
For every $N\ge0$ and $\varepsilon>0$ there is a chain with $n=2N+4$ states, an induced chain and a set $A$ with two
elements such that $\hE_{\hat x}\widehat{T}_{\hA}\ge(N+1-\varepsilon)\max_y\E_yT_A$ for some $\hat x$.
\end{proposition}

Section~\ref{sec:pot} collects identities for hitting times, Theorem~\ref{thm:merge} is proved in
Section~\ref{sec:merge}, Lemma~\ref{lem:relax} and the upper bound of Theorem~\ref{thm:main} in
Section~\ref{sec:upper}, the lower bound in Section~\ref{sec:comb} and Proposition~\ref{prop:hit} in
Section~\ref{sec:hit}.

\section{Hitting-time potentials}\label{sec:pot}
Throughout, $H(x,j)=\E_xT_j$ and $F_{ij}=\E_iT_j-\E_\pi T_j$, so that $\tmix=\max_{i,j}F_{ij}$ by~\eqref{eq:lw}.
Functions on $I$ are column vectors; $\omega$ is the vector $(w_x)_{x\in I}$, $\one$ the constant vector $1$ and
$\delta_x$ the indicator of $x$. The Laplacian $L$ has entries $L_{xy}=-w_{xy}$ for $x\ne y$ and
$L_{xx}=\sum_{y\ne x}w_{xy}$, so that loops do not enter it, and the Dirichlet form is
\[
\mathcal E(g,g')=g^{\mathsf T}Lg'=\frac12\sum_{x,y}w_{xy}\,(g_x-g_y)(g'_x-g'_y).
\]
The kernel of $L$ consists of the constants, and $\one^{\mathsf T}Lg=0$ for every $g$. For a state $y$ and a set $Z$
not containing it, $R(y\lra Z)$ is the effective resistance between $y$ and $Z$ in the network with conductances
$w_{xy}$, $x\ne y$. For two states $u\ne v$ the function $h_x=\Prob_x(T_u<T_v)$ satisfies~\cite[Section~2.2]{LP}
\begin{equation}\label{eq:harm}
Lh=\frac{\delta_u-\delta_v}{R(u\lra v)},\qquad \mathcal E(h,h)=\frac1{R(u\lra v)} .
\end{equation}

Fix a state $o$. Since the network is connected, the matrix obtained from $L$ by deleting the row and the column of
$o$ is invertible; let $\Gamma$ be its inverse, extended by zeros in the row and the column of $o$. Then $\Gamma$ is
symmetric and $L\Gamma\delta_j=\delta_j-\delta_o$ for $j\ne o$, so that $\Gamma\delta_j$ is the voltage, vanishing at
$o$, of the unit current flow from $j$ to $o$. Put $\zeta=\Gamma\omega$ and $\bar\zeta=\sum_x\pi_x\zeta_x$.

\begin{lemma}\label{lem:pot}
\begin{enumerate}
\item[(i)] $H(i,j)=\zeta_i-\zeta_j-w(\Gamma_{ij}-\Gamma_{jj})$ for all $i,j$; in particular $\zeta_i=H(i,o)$.
\item[(ii)] $F_{ij}=\zeta_i+\zeta_j-\bar\zeta-w\Gamma_{ij}$. In particular $F$ is symmetric,
$\sum_x\pi_xF_{ix}=0$, and $LF\delta_j=\omega-w\delta_j$.
\item[(iii)] For $i\in I$ let $\kappa_i=\max_jF_{ij}$ and let $\varphi_i$ be the function
$\varphi_i(x)=\kappa_i-F_{ix}$. Then $\kappa_i\ge0$, $\varphi_i\ge0$, $\sum_x\pi_x\varphi_i(x)=\kappa_i$,
$\tmix=\max_i\kappa_i$, and for every function $g$ on $I$
\begin{equation}\label{eq:energy}
\mathcal E(\varphi_i,g)=w\Bigl(g_i-\sum_x\pi_xg_x\Bigr),\qquad \mathcal E(\varphi_i,\varphi_j)=-wF_{ij}.
\end{equation}
\end{enumerate}
\end{lemma}

\begin{proof}
(i) For $x\ne j$, conditioning on the first step gives $w_xH(x,j)=w_x+\sum_yw_{xy}H(y,j)$, that is,
$(LH(\cdot,j))_x=w_x$; as the entries of $LH(\cdot,j)$ sum to zero, $(LH(\cdot,j))_j=w_j-w$. Thus
$LH(\cdot,j)=\omega-w\delta_j$ and $H(j,j)=0$, and these two conditions determine $H(\cdot,j)$. Now
$L\zeta=\sum_{j\ne o}w_j(\delta_j-\delta_o)=\omega-w\delta_o$, so the right side of (i), as a function of $i$, has
Laplacian $\omega-w\delta_o-w(\delta_j-\delta_o)=\omega-w\delta_j$ when $j\ne o$ and $\omega-w\delta_o$ when $j=o$,
and it vanishes at $i=j$. This proves the formula, and $\zeta_o=0$, $\Gamma_{io}=0$ give $\zeta_i=H(i,o)$.

(ii) By the symmetry of $\Gamma$, $\sum_i\pi_i\Gamma_{ij}=\zeta_j/w$. Averaging (i) over $i$ with weights $\pi_i$
gives $\E_\pi T_j=\bar\zeta-2\zeta_j+w\Gamma_{jj}$, and subtracting this from (i) gives the formula for $F_{ij}$,
which is symmetric. Each column of $F$ has $\pi$-mean zero by definition, hence so does each row, and
$LF\delta_j=LH(\cdot,j)$.

(iii) A row of $F$ has $\pi$-mean zero, so its maximum $\kappa_i$ is nonnegative, and $\varphi_i\ge0$ has
$\pi$-mean $\kappa_i$. The identity $\tmix=\max_i\kappa_i$ is~\eqref{eq:lw}. By symmetry of $F$ and (ii),
$L\varphi_i=-LF\delta_i=w\delta_i-\omega$, so $\mathcal E(\varphi_i,g)=g^{\mathsf T}L\varphi_i=wg_i-\sum_xw_xg_x$.
With $g=\varphi_j$ this is $w(\kappa_j-F_{ji}-\kappa_j)=-wF_{ij}$.
\end{proof}

We also use two facts about networks whose graph, loops removed, is a tree.

\begin{lemma}\label{lem:bridge}
Let $x\ne y$ be states with $w_{xy}>0$ such that every path from $x$ to $y$ in the network uses the edge $\{x,y\}$,
and let $V_{xy}$ be the set of states that can be reached from $x$ without visiting $y$. Then
$\E_xT_y=w(V_{xy})/w_{xy}$, where $w(V_{xy})=\sum_{z\in V_{xy}}w_z$.
\end{lemma}
\begin{proof}
Up to time $T_y$ the walk started at $x$ stays in $V_{xy}$, so $\E_xT_y$ is the same for the network on
$V_{xy}\cup\{y\}$ in which $y$ keeps only the edge to $x$. There the stationary probability of $y$ is
$w_{xy}/(w(V_{xy})+w_{xy})$, so the mean return time to $y$ is $1+w(V_{xy})/w_{xy}$, and it also equals
$1+\E_xT_y$.
\end{proof}

\begin{lemma}\label{lem:tree}
Suppose that the graph of the network, loops removed, is a tree. For a state $x$ let $[x,o]$ be the set of edges of
the path from $x$ to $o$, for an edge $e$ let $w_e$ be its weight, and let $S_e$ be the total weight $\sum_zw_z$ of
the states separated from $o$ by $e$. Then
\[
\Gamma_{ij}=\sum_{e\in[i,o]\cap[j,o]}\frac1{w_e},\qquad H(x,o)=\sum_{e\in[x,o]}\frac{S_e}{w_e},\qquad
\E_\pi T_o=\frac1w\sum_e\frac{S_e^2}{w_e},
\]
the last sum being over all edges of the tree.
\end{lemma}
\begin{proof}
On a tree the unit current flow from $j$ to $o$ runs along the path from $j$ to $o$, and its voltage at $i$ is the
sum of the resistances $1/w_e$ of the edges of that path through which the path from $i$ to $o$ also passes. The walk
from $x$ to $o$ visits every state of the path between them in turn, and by Lemma~\ref{lem:bridge} the edge $e$ of
that path is crossed towards $o$ in mean time $S_e/w_e$. Averaging over $x$ with weights $\pi_x$, the edge $e$ is
counted with total weight $S_e/w$.
\end{proof}

\section{Lumping two states}\label{sec:merge}
Let $u\ne v$ be states, let $f$ send $u$ and $v$ to one state $\hat v$, send no other state to $\hat v$ and be injective on
$J=I\setminus\{u,v\}$, and
let $\hF$ be the matrix $F$ of $\hX$. Let $h_x=\Prob_x(T_u<T_v)$, $R=R(u\lra v)$ and $\bar h=\sum_x\pi_xh_x$, and for
$i\in I$ put
\[
\beta_i=\varphi_i(u)-\varphi_i(v)=F_{iv}-F_{iu}.
\]

\begin{proposition}\label{prop:rankone}
For all $i,j\in I$,
\[
\beta_i=wR\,(h_i-\bar h),\qquad \hF_{f(i)f(j)}=F_{ij}+\frac{\beta_i\beta_j}{wR},
\]
and
\begin{equation}\label{eq:twoterms}
\hF_{f(i)f(j)}\ \le\ \kappa_i+\kappa_j .
\end{equation}
\end{proposition}

\begin{proof}
Take $o=v$ in Section~\ref{sec:pot}. By~\eqref{eq:harm} the voltage of the unit current flow from $u$ to $v$ is $Rh$,
so $\Gamma\delta_u=Rh$; in particular $\Gamma_{uu}=R$, $\zeta_u=\sum_x\Gamma_{ux}w_x=wR\bar h$ and $\zeta_v=0$.
Lemma~\ref{lem:pot}(ii) gives $\beta_i=\zeta_v-\zeta_u+w(\Gamma_{iu}-\Gamma_{iv})=wR(h_i-\bar h)$.

Let $\widetilde\Gamma=\Gamma-Rhh^{\mathsf T}$. Its rows and columns at $u$ and $v$ vanish. The rows of
$L\Gamma$ indexed by $J$, restricted to the columns indexed by $J$ and by $u$, give
$L_{JJ}\Gamma_{JJ}+L_{Ju}\Gamma_{uJ}=I_J$ and $L_{JJ}\Gamma_{Ju}+L_{Ju}\Gamma_{uu}=0$, where the subscripts denote
submatrices; hence $L_{JJ}(\Gamma_{JJ}-\Gamma_{Ju}\Gamma_{uJ}/\Gamma_{uu})=I_J$, that is,
$\widetilde\Gamma_{JJ}=L_{JJ}^{-1}$. The Laplacian of $\hX$ with the row and the column of $\hat v$ deleted is
$L_{JJ}$, the total weight of $\hX$ is $w$, and $\hpi_{f(x)}=\pi_x$ for $x\in J$, $\hpi_{\hat v}=\pi_u+\pi_v$. So the
matrix $\Gamma$ of $\hX$ grounded at $\hat v$ has entries $\widetilde\Gamma_{ij}$ at $(f(i),f(j))$.

For a matrix $A$ indexed by $I$ write $\mathcal C(A)_{ij}=A_{ij}-(A\pi)_i-(\pi^{\mathsf T}A)_j+\pi^{\mathsf T}A\pi$.
Lemma~\ref{lem:pot}(ii) says that $F=-w\,\mathcal C(\Gamma)$. If $A_{ij}$ depends only on $f(i)$ and $f(j)$, then
$\mathcal C(A)_{ij}$ equals the same expression formed with the quotient matrix and $\hpi$ at $(f(i),f(j))$. Applied
to $\widetilde\Gamma$, and since $\mathcal C(hh^{\mathsf T})_{ij}=(h_i-\bar h)(h_j-\bar h)$,
\[
\hF_{f(i)f(j)}=-w\,\mathcal C(\widetilde\Gamma)_{ij}=F_{ij}+wR\,(h_i-\bar h)(h_j-\bar h)
=F_{ij}+\frac{\beta_i\beta_j}{wR}.
\]

If $\beta_i\beta_j\le0$ this gives $\hF_{f(i)f(j)}\le F_{ij}\le\kappa_i\le\kappa_i+\kappa_j$. Suppose that
$\beta_i\beta_j>0$. Exchanging the names of $u$ and $v$ changes the signs of $\beta_i$ and $\beta_j$ together and
leaves the statement unchanged, so we may assume $\beta_i,\beta_j>0$. Put $g_i=\varphi_i/\beta_i$,
$g_j=\varphi_j/\beta_j$ and let $\eta=\min(g_i,g_j)$ pointwise. These functions are nonnegative, and
$g_i(u)-g_i(v)=g_j(u)-g_j(v)=1$, hence $\eta(u)-\eta(v)=1$. The function $\eta-\eta(v)\one-h$ vanishes at $u$ and at
$v$, so by~\eqref{eq:harm} it is $\mathcal E$-orthogonal to $h$, and
\[
\mathcal E(\eta,\eta)\ \ge\ \mathcal E(h,h)=\frac1R .
\]
Let $\xi_i=g_i-\eta$ and $\xi_j=g_j-\eta$. They are nonnegative and at every state at least one of them vanishes, so
$\mathcal E(\xi_i,\xi_j)=-\sum_{x\ne y}w_{xy}\xi_i(x)\xi_j(y)\le0$. Expanding $g_i=\eta+\xi_i$ and $g_j=\eta+\xi_j$,
\[
\mathcal E(g_i,\eta)+\mathcal E(g_j,\eta)-\mathcal E(g_i,g_j)=\mathcal E(\eta,\eta)-\mathcal E(\xi_i,\xi_j)
\ \ge\ \frac1R .
\]
Let $\bar\eta=\sum_x\pi_x\eta(x)\ge0$. By~\eqref{eq:energy} the left side equals
\[
\frac w{\beta_i}\bigl(\eta(i)-\bar\eta\bigr)+\frac w{\beta_j}\bigl(\eta(j)-\bar\eta\bigr)+\frac{wF_{ij}}{\beta_i\beta_j},
\]
and $\eta(i)\le g_j(i)=(\kappa_j-F_{ij})/\beta_j$, $\eta(j)\le g_i(j)=(\kappa_i-F_{ij})/\beta_i$. Hence
\[
\frac1R\ \le\ \frac w{\beta_i\beta_j}\bigl(\kappa_i+\kappa_j-F_{ij}-(\beta_i+\beta_j)\bar\eta\bigr),
\]
and multiplying by $\beta_i\beta_j/w$ and using the formula for $\hF$,
\begin{equation}\label{eq:strict}
\hF_{f(i)f(j)}\ \le\ \kappa_i+\kappa_j-(\beta_i+\beta_j)\,\bar\eta . \qedhere
\end{equation}
\end{proof}

\begin{proof}[Proof of Theorem~\ref{thm:merge}]
By~\eqref{eq:lw} and Lemma~\ref{lem:pot}(iii), $\tmix(\hX)$ is the largest entry of $\hF$, and every entry is
$\hF_{f(i)f(j)}$ for some $i,j\in I$. If $\beta_i\beta_j\le0$ that entry is at most $\kappa_i\le\tmix(X)$. Otherwise,
in the notation of the proof above, $\eta(u)=\eta(v)+1\ge1$, so $\bar\eta\ge\pi_u>0$, and~\eqref{eq:strict} gives
$\hF_{f(i)f(j)}<\kappa_i+\kappa_j\le2\tmix(X)$. Finally $\tmix(X)>0$: otherwise every row of $F$, having
$\pi$-mean zero, would vanish, against $LF\delta_j=\omega-w\delta_j\ne0$. As there are finitely many entries,
$\tmix(\hX)<2\tmix(X)$. That the constant $2$ is optimal is shown by the comb of Section~\ref{sec:comb} with $N=0$
(Proposition~\ref{prop:comb}).
\end{proof}

\section{The upper bound}\label{sec:upper}

\begin{proof}[Proof of Lemma~\ref{lem:relax}]
Replacing $g$ by a multiple we may assume that $g(i)=1=\max_x|g(x)|$. Since $Pg=\lambda g$, where $P=(p_{xy})$, the
process
\[
\mathcal M_t=g(X_t)+(1-\lambda)\sum_{s<t}g(X_s),\qquad t\ge0,
\]
is a martingale with $\mathcal M_0=1$ under $\Prob_i$, with respect to the filtration generated by the chain and by
the randomization used by $U$. We may assume that $\E_iU<\infty$. The increments of $\mathcal M$ are bounded by $2$,
so the optional stopping theorem gives $\E_i\mathcal M_U=1$. As $\lambda\ne1$,
$\sum_x\pi_xg(x)=\sum_x\pi_x(Pg)(x)=\lambda\sum_x\pi_xg(x)$ forces $\sum_x\pi_xg(x)=0$. Hence $\E_ig(X_U)=0$ and
\[
1=(1-\lambda)\,\E_i\sum_{s<U}g(X_s)\le(1-\lambda)\,\E_iU .
\]
For a reversible chain all eigenvalues are real; taking $\lambda=\lambda_2$ gives $\tmix\ge\min_U\E_iU\ge\tau_2$.
\end{proof}

\begin{proof}[Proof of the upper bound in Theorem~\ref{thm:main}]
Let $X$ have $m\le n$ states and let $\hX$ be an induced chain with $q$ states. If $q=m$ the map $f$ is a bijection
and $\hX$ is $X$ with its states renamed, and if $q=1$ then $\tmix(\hX)=0$. Otherwise $m\ge3$ and $2\le q\le m-1$.
The lumping is a composition of $m-q$ lumpings of two states, so Theorem~\ref{thm:merge} gives
$\tmix(\hX)\le2^{m-q}\tmix(X)$. We also have
\begin{equation}\label{eq:chain}
\tmix(\hX)\ \le\ \tau_0(\hX)\ \le\ (q-1)\,\tau_2(\hX)\ \le\ (q-1)\,\tau_2(X)\ \le\ (q-1)\,\tmix(X).
\end{equation}
For the first inequality, choose a state $J$ of $\hX$ with distribution $\hpi$, independently of the chain, and stop
at the hitting time of $J$; the chain is then in the stationary distribution, and the mean of this stopping time is
$\tau_0(\hX)$ from every starting state~\cite[Section~4.3]{AF}. The second is the eigentime identity
$\tau_0=\sum_{k\ge2}(1-\lambda_k)^{-1}$~\cite[Proposition~3.13 and Lemma~4.24]{AF} for the reversible chain $\hX$,
each of whose $q-1$ terms is at most $\tau_2(\hX)$. The third is part of the contraction
principle~\cite[Proposition~4.44]{AF}: a function $g$ on $\hI$ and the function $g\circ f$ on $I$ have the same mean,
variance and Dirichlet form under $\hpi$ and $\pi$, so the spectral gap of $\hX$ is at least that of $X$. The fourth
is Lemma~\ref{lem:relax}. Hence the ratio is at most $\min(2^{m-q},q-1)\le U(m)$. Finally $U(m)\le U(n)$, because
$2^{m-q}\le2^{n-q}$, and $U(n)\le\max(1,n-2)$, because each term is at most $q-1\le n-2$.
\end{proof}

\begin{remark}\label{rem:pimin}
Write $\pi_*=\min_x\pi_x$. Then $\tmix\le\frac{4e}{e-1}\,\tau_2\,(1+\frac12\log(1/\pi_*))$: in
continuous time this combines $\tmix\le\frac{e}{e-1}\tau_1^{(1)}$, $\tau_1^{(1)}\le4\tau_1$ and
$\tau_1\le\tau_2(1+\frac12\log(1/\pi_*))$ \cite[Lemmas~4.12, 4.11 and~4.23]{AF}, where $\tau_1$ and $\tau_1^{(1)}$ are the first times at which the
largest total variation distance between the distributions from two starting states and the separation from $\pi$
fall to $e^{-1}$~\cite[Section~4.3]{AF}, and $\tmix$ of the discrete-time chain
is at most that of its continuization~\cite[Lemma~4.16]{AF}. Since $\tau_2(\hX)\le\tau_2(X)$ and $\min_{\hat x}\hpi_{\hat x}\ge\pi_*$,
Lemma~\ref{lem:relax} gives
\[
\tmix(\hX)\ \le\ \frac{4e}{e-1}\Bigl(1+\frac12\log\frac1{\pi_*}\Bigr)\tmix(X),
\]
which is the form taken for $\tmix$ by the bound used by Hermon and Peres in~\cite[Remark~1.8]{HP}. Together with
Theorem~\ref{thm:main}, the factor is at most $\min\bigl(U(n),\ \frac{4e}{e-1}(1+\frac12\log(1/\pi_*))\bigr)$. For the
simple random walk on a graph with bounded degrees $\log(1/\pi_*)$ is of order $\log n$, and the examples
of~\cite{HP} show that this order is attained. In the comb of Section~\ref{sec:comb}, $\log(1/\pi_*)$ is of order
$(N+1)\log t$.
\end{remark}

\section{The comb}\label{sec:comb}
Fix an integer $N\ge0$ and a real number $t\ge2$, and put $r=t^3$ and $\alpha=t$. The \emph{comb} is the network on
the $2N+4$ states $a$, $b$, $c$, $z_0,\dots,z_N$, $y_0,\dots,y_{N-1}$ with the edge weights
\begin{gather*}
w_{az_0}=1,\qquad w_{bz_N}=\alpha r^{N+1},\qquad w_{cz_N}=r^{N+2},\\
w_{z_kz_{k+1}}=r^{k+1},\qquad w_{y_kz_{k+1}}=\alpha r^{k+1}\qquad(0\le k<N),
\end{gather*}
the loop weights
\begin{gather*}
w_{aa}=t-1,\qquad w_{bb}=(t-1)\alpha r^{N+1},\qquad w_{cc}=\ell=(N+2)\,r^{N+4},\\
w_{y_ky_k}=(t-1)\alpha r^{k+1}\qquad(0\le k<N),
\end{gather*}
and all other weights zero. Without its loops the comb is a tree: the spine $a,z_0,\dots,z_N,c$, the teeth $y_k$
hanging from $z_{k+1}$, and the leaf $b$ hanging from $z_N$. From $a$, from $b$ and from each tooth the walk leaves
with probability $1/t$ at each step. The map $f$ sends $z_k$ and $y_k$ to one state $\hat z_k$ for $0\le k<N$, sends
$z_N$ and $c$ to one state $\hat c$, and keeps $a$ and $b$. So $\hX$ has $N+3$ states and is obtained by $N+1$
lumpings of two states. For $N=0$ the comb is a star with centre $z_0$, and only $z_0$ and $c$ are lumped.

Put $\sigma_k=\sum_{l=1}^kr^l$ and $B=2+\alpha(t+1)$, and
\begin{gather*}
S_k=t+1+r^{k+1}+B\sigma_k,\qquad \hS_k=S_k+t\alpha r^{k+1}\qquad(0\le k<N),\\
S_N=t+1+r^{N+2}+B\sigma_N+(t+1)\alpha r^{N+1},\\
M=t+\sum_{k=0}^{N-1}\frac{S_k}{r^{k+1}}+\frac{S_N}{r^{N+2}},\qquad
\hM=t+\sum_{k=0}^{N-1}\frac{\hS_k}{(1+\alpha)r^{k+1}} .
\end{gather*}

\begin{proposition}\label{prop:comb}
For the comb $X$ and its induced chain $\hX$:
\begin{enumerate}
\item[(i)] $S_k$ is the total weight of the states separated from $c$ by the edge $\{z_k,z_{k+1}\}$ ($k<N$), $S_N$ is
that of the states other than $c$, and $w=\ell+r^{N+2}+S_N$. Moreover $M=\E_aT_c=\max_x\E_xT_c$,
$\hM=\hE_a\widehat T_{\hat c}$, $\hE_b\widehat T_{\hat c}=t$, and $\hE_{\hat x}\widehat T_{\hat c}\le\hM$ for
$\hat x\ne b$.
\item[(ii)] If $w\ge Mr^{N+2}$, then $\tmix(X)=M-\E_\pi T_c$. If moreover $N=0$ or $w\ge(\hM-t)(1+\alpha)r^N$, then
$\tmix(\hX)=\hM+t-\hE_{\hpi}\widehat T_{\hat c}$.
\item[(iii)] As $t\to\infty$ with $N$ fixed, the conditions of (ii) hold, $M=t+N+1+O(1/t)$,
$\hM=(N+1)t+O(1)$, and $\E_\pi T_c$ and $\hE_{\hpi}\widehat T_{\hat c}$ tend to zero. Hence
\[
\frac{\tmix(\hX)}{\tmix(X)}\to N+2,\qquad \frac{\hM}{M}\to N+1 .
\]
\end{enumerate}
\end{proposition}

\begin{proof}
(i) The state weights are $w_a=t$, $w_{y_k}=t\alpha r^{k+1}$, $w_b=t\alpha r^{N+1}$, $w_c=r^{N+2}+\ell$, and for the
spine $w_{z_0}=1+r$ and $w_{z_k}=r^k+r^{k+1}+\alpha r^k$ ($1\le k<N$), $w_{z_N}=r^N+\alpha r^N+\alpha r^{N+1}+r^{N+2}$
when $N\ge1$, while $w_{z_0}=1+\alpha r+r^2$ when $N=0$. The states separated from $c$ by $\{z_k,z_{k+1}\}$ are $a$,
$z_0,\dots,z_k$ and $y_0,\dots,y_{k-1}$. Their total weight is $t+1+r$ for $k=0$, and passing from $k-1$ to $k$ adds
$w_{z_k}+w_{y_{k-1}}=r^{k+1}+(B-1)r^k$, which is also the increase of $t+1+r^{k+1}+B\sigma_k$. The same count with
$w_b$ added gives $S_N$, and $w=w_c+S_N$.

By Lemma~\ref{lem:bridge}, $\E_aT_{z_0}=t$, $\E_{z_k}T_{z_{k+1}}=S_k/r^{k+1}$, $\E_{z_N}T_c=S_N/r^{N+2}$,
$\E_{y_k}T_{z_{k+1}}=t$ and $\E_bT_{z_N}=t$. The walk from $x$ to $c$ visits every state of the path between them, so
$\E_xT_c$ is the sum of these terms along the path. From $a$ this sum is $M$. From any other state it is either a sum
of some of the spine terms, or $t$ followed by some of the spine terms, and in both cases at most $M$.

The induced chain is again a tree: the path $a,\hat z_0,\dots,\hat z_{N-1},\hat c$ with the leaf $b$ attached to
$\hat c$. Writing $\hat z_N=\hat c$, the edge $\{\hat z_k,\hat z_{k+1}\}$ has weight
$w_{z_kz_{k+1}}+w_{y_kz_{k+1}}=(1+\alpha)r^{k+1}$, and the states separated from $\hat c$ by it are $a,\hat z_0,\dots,
\hat z_k$, of total weight $S_k+w_{y_k}=\hS_k$. The edge $\{a,\hat z_0\}$ (or $\{a,\hat c\}$ if $N=0$) has weight $1$
and $\{b,\hat c\}$ has weight $\alpha r^{N+1}$. Lemma~\ref{lem:bridge} and the argument above give the assertions
about $\hX$.

(ii) Take $o=c$ in Section~\ref{sec:pot}. By Lemma~\ref{lem:pot}, $F_{ij}=H(i,c)+H(j,c)-\E_\pi T_c-w\Gamma_{ij}$. For
$j=c$ this is $H(i,c)-\E_\pi T_c\le M-\E_\pi T_c$, with equality at $i=a$. For $i,j\ne c$ the edge $\{z_N,c\}$ lies in
$[i,c]\cap[j,c]$, so $\Gamma_{ij}\ge r^{-(N+2)}$ by Lemma~\ref{lem:tree}, and
$F_{ij}\le2M-\E_\pi T_c-wr^{-(N+2)}\le M-\E_\pi T_c$. By~\eqref{eq:lw}, $\tmix(X)=M-\E_\pi T_c$.

In $\hX$ take $o=\hat c$ and write $\widehat H$, $\widehat\Gamma$ for the corresponding quantities. The paths from $a$
and from $b$ to $\hat c$ share no edge, so $\hF_{ab}=\hM+t-\hE_{\hpi}\widehat T_{\hat c}$. Every other entry is at most
this: entries with $\hat c$ are at most $\hM-\hE_{\hpi}\widehat T_{\hat c}$; entries $\hF_{b\hat x}$ with $\hat x\ne
b$ are at most $\widehat H(\hat x,\hat c)+t-\hE_{\hpi}\widehat T_{\hat c}$; $\hF_{bb}\le2t-\hE_{\hpi}\widehat
T_{\hat c}$ and $t\le\hM$; if $N=0$ the only remaining entry is $\hF_{aa}\le2t-\hE_{\hpi}\widehat T_{\hat c}$ and
$\hM=t$; and if $N\ge1$, two states of the path other than $\hat c$ share the edge $\{\hat z_{N-1},\hat c\}$, so
$\widehat\Gamma\ge1/((1+\alpha)r^N)$ for them and their entry is at most
$2\hM-\hE_{\hpi}\widehat T_{\hat c}-w/((1+\alpha)r^N)\le\hM+t-\hE_{\hpi}\widehat T_{\hat c}$.

(iii) Since $\sigma_k\le2r^k$, $B=t^2+t+2$ and $r=t^3$, each of $S_k/r^{k+1}$ ($k<N$) and $S_N/r^{N+2}$ is
$1+O(1/t)$, so $M=t+N+1+O(1/t)$, and $\hS_k/((1+\alpha)r^{k+1})=(t^2+1+O(1/t))/(t+1)=t-1+O(1/t)$, so
$\hM=(N+1)t-N+O(1/t)$. As $S_N=O(r^{N+2})$ and $w\ge\ell=(N+2)r^{N+4}$, the conditions of (ii) hold for large $t$.
Next, $\E_\pi T_c\le(1-\pi_c)M=(S_N/w)M=O(t^{-5})$, and since the states of $\hX$ other than $\hat c$ carry
stationary mass at most $S_N/w$, also $\hE_{\hpi}\widehat T_{\hat c}\le(S_N/w)\hM=O(t^{-5})$. By (ii) the ratio is
$((N+2)t-N+o(1))/(t+N+1+o(1))\to N+2$.
\end{proof}

For $r\gg\alpha t$ a walk on the spine of the comb rarely enters a tooth, and each tooth or leaf delays a walk that
starts there by its holding time, of mean $t$. After lumping, each state $\hat z_k$ inherits the holding loop of the
tooth $y_k$, and since $\alpha$ is large the walk from $a$ to $\hat c$ is delayed by about $t$ at each of the $N$
levels. In $X$ the far ends $a$ and $b$ are each about $t$ away from $c$ and the heavy state $c$ separates them; in
$\hX$ the delays add up along the spine, and the leaf $b$ adds one more $t$.

\begin{proof}[Proof of the lower bound in Theorem~\ref{thm:main}]
The identity map gives $K(n)\ge1$. For $n\ge4$ let $N=\lfloor n/2\rfloor-2$. The comb has $2N+4\le n$ states, and by
Proposition~\ref{prop:comb} the ratio $\tmix(\hX)/\tmix(X)$ tends to $N+2=\lfloor n/2\rfloor$ as $t\to\infty$. The
stated values of $K(n)$ and the bounds on $K(7)$ follow from Table~\ref{tab:bounds}.
\end{proof}

Table~\ref{tab:comb} gives the values for $t=1000$, computed in rational arithmetic from~\eqref{eq:lw} and rounded;
the closed forms of Proposition~\ref{prop:comb} agree with them exactly, and the conditions of
Proposition~\ref{prop:comb}(ii) hold already for $t=2$ when $N\le3$.

\begin{table}[tbp]
\centering
\begin{tabular}{rrrrrr}
\toprule
$N$ & states & $\tmix(X)$ & $\tmix(\hX)$ & $\tmix(\hX)/\tmix(X)$ & $\hM/M$\\
\midrule
0 & 4 & 1001.001 & 2000.000 & 1.9980 & 0.9990\\
1 & 6 & 1002.001 & 2999.002 & 2.9930 & 1.9950\\
2 & 8 & 1003.002 & 3998.004 & 3.9860 & 2.9890\\
3 & 10 & 1004.003 & 4997.006 & 4.9771 & 3.9811\\
\bottomrule
\end{tabular}
\caption{The comb with $t=1000$ and its induced chain.}\label{tab:comb}
\end{table}

\section{Hitting times of sets}\label{sec:hit}
For a non-empty proper subset $Z$ of $I$ and $x,y\notin Z$ let $G_Z(x,y)$ be the expected number of visits to $y$
before $T_Z$ for the walk started at $x$. Then~\cite[Section~2.2]{LP}
\begin{equation}\label{eq:green}
\E_xT_Z=\sum_{y\notin Z}G_Z(x,y),\qquad G_Z(x,y)=\Prob_x(T_y<T_Z)\,G_Z(y,y),\qquad G_Z(y,y)=w_y\,R(y\lra Z).
\end{equation}
The formula for $G_Z(y,y)$ holds in the presence of loops: the walk at $y$ reaches $Z$ before returning to $y$ with
probability $1/(w_yR(y\lra Z))$, and a loop at $y$ changes $w_y$ only.

\begin{proof}[Proof of Proposition~\ref{prop:hit}]
The network of $\hX$ is obtained from that of $X$ by gluing together the states in each set $f^{-1}(\hat y)$; edges
inside such a set become loops. By Rayleigh's monotonicity principle~\cite[Section~2.4]{LP} gluing does not increase
effective resistances, so $\widehat{R}(\hat y\lra\hA)\le R(y\lra A)$ for every $\hat y\in\hI\setminus\hA$ and every
$y\in f^{-1}(\hat y)$, where $\widehat R$ is the effective resistance in the network of $\hX$. With~\eqref{eq:green}
applied to both chains,
\begin{align*}
\hE_{\hat x}\widehat{T}_{\hA}\ &\le\ \sum_{\hat y\notin\hA}\hw_{\hat y}\,\widehat{R}(\hat y\lra\hA)
=\sum_{\hat y\notin\hA}\ \sum_{y\in f^{-1}(\hat y)}w_y\,\widehat{R}(\hat y\lra\hA)\\
&\le\ \sum_{y\notin A}w_y\,R(y\lra A)=\sum_{y\notin A}G_A(y,y)\ \le\ \sum_{y\notin A}\E_yT_A .
\end{align*}
For the second statement take the comb of Section~\ref{sec:comb}, $\hA=\{\hat c\}$, $A=\{z_N,c\}$ and $\hat x=a$.
Then $\max_y\E_yT_A\le\max_y\E_yT_c=M$ and $\hE_a\widehat T_{\hat c}=\hM$, and $\hM/M\to N+1$ by
Proposition~\ref{prop:comb}(iii).
\end{proof}

The program \texttt{merge\_and\_comb\_checks.py} in the repository named below checks, in rational arithmetic, the
formula for $\hF$ and inequality~\eqref{eq:twoterms} for every pair of states and every lumping of two states in $100$
random networks with three to six states ($850$ lumpings), and the closed forms of Proposition~\ref{prop:comb} for
$N\le3$ and $t\in\{2,3,10,100,1000\}$. The program \texttt{lumping\_checks.py} checks Proposition~\ref{prop:hit},
the inequalities $\tmix(\hX)\le\tau_0(\hX)\le\tau_0(X)\le(m-1)\,\tmix(X)$ and the upper bound of
Theorem~\ref{thm:main} on $400$ random networks with three to six states and random partitions.

\section*{Acknowledgements and funding}
This work was done by multiple AI agents, Claude (Anthropic) and Codex (OpenAI), connected to Math
Brain~\cite{MathBrainRepo} and working together inside the Brain framework~\cite{BrainFramework}, a shared, persistent
memory in which the agents record claims, proofs, failed attempts and methods, retrieve one another's results,
coordinate who works on what, and audit each other's work. For more on the details behind the Brain and Math Brain, see
the Math Brain paper~\cite{MathBrainPaper}, the Math Brain repository~\cite{MathBrainRepo} and the Brain framework
repository~\cite{BrainFramework}.

The research presented in this paper was supported by the European Research Council (ERC) under the European Union's
Horizon Europe research and innovation programme (grant agreement No.~101041711), by the Simons Foundation as part of
the Collaboration on the Mathematical and Scientific Foundations of Deep Learning, by Heights Labs, by Convex Nexus
Capital, by the Israel Science Foundation (grant number 2258/19), and by the Israel Science Foundation (ISF Grant
4101/25). Computational resources were provided in part by Google Cloud through its education credits program (credit
530550103).

\section*{Declaration of competing interest}
The author holds equity in Heights Labs and in Convex Nexus Capital and has a management role at Convex Nexus Capital.
These entities are included in the funding declaration above.

\section*{Declaration of generative AI and AI-assisted technologies}
The paper comes out of Math Brain, an ongoing project on the use of AI systems for research mathematics, and its
results were developed with the Brain~\cite{BrainFramework} and Math Brain~\cite{MathBrainRepo}. In October 2026 an
AI agent connected to the Brain, working from the project's list of open problems in probability, found examples in
which lumping increases $\tmix$, among them the example of Remark~\ref{rem:four}. In November 2026 a Claude Fable
agent (Anthropic) obtained Theorem~\ref{thm:main}, Lemma~\ref{lem:relax} and Propositions~\ref{prop:hit}
and~\ref{prop:comb} with their proofs, wrote and ran the programs and prepared the manuscript. A Codex agent (OpenAI) then read the proofs and recomputed
the examples with programs of its own. These checks are not peer review. The author directed the work and is responsible for the content of the paper and for decisions
concerning its dissemination.

\section*{Code and data availability}
Programs and outputs:\\
\url{https://github.com/yspennstate/lumping-mixing-time}.

The paper is part of the Math Brain project, whose public repository is
\url{https://github.com/yspennstate/math-brain-public}~\cite{MathBrainRepo}.


\appendix
\section{Notation}\label{app:notation}
Symbols are listed in the order in which they first appear. Letters used only inside one proof are not listed.

\medskip
\begin{center}
\small
\begin{tabular}{@{}lp{0.6\textwidth}l@{}}
\toprule
Symbol & Meaning & First appears\\
\midrule
$I$, $n$ & the state space of $X$ and its number of elements & Section~\ref{sec:intro}\\
$w_{xy}$, $w_x$, $w$ & the weight of $\{x,y\}$, the weight of a state, the total weight & Section~\ref{sec:intro}\\
$X$, $p_{xy}$, $\pi$ & the random walk, its transition probabilities, its stationary law & Section~\ref{sec:intro}\\
$T_y$, $T_A$ & the hitting time of a state $y$ and of a set $A$ & Section~\ref{sec:intro}\\
$\E_\mu$, $\Prob_\mu$ & expectation and probability for the chain started from $\mu$ & Section~\ref{sec:intro}\\
$f$, $\hI$ & the lumping map and its image set & Section~\ref{sec:intro}\\
$\hX$, $\hw$, $\hpi$ & the induced chain, its weights, its stationary law & Section~\ref{sec:intro}\\
$\hE$, $\widehat T$ & expectation and hitting times for $\hX$ & Section~\ref{sec:intro}\\
$\tmix$ & the stopping-rule mixing time, given by~\eqref{eq:lw} & Section~\ref{sec:intro}\\
$\tau^*$, $\tau_0$ & the maximal mean commute time and the average hitting time & Section~\ref{sec:intro}\\
$\lambda_k$, $\tau_2$ & the $k$-th largest eigenvalue of $(p_{xy})$ and $1/(1-\lambda_2)$ & Section~\ref{sec:intro}\\
$\tau_c$ & the Cheeger time constant & Section~\ref{sec:intro}\\
$K(n)$ & the largest factor by which lumping can increase $\tmix$ on at most $n$ states & Section~\ref{sec:intro}\\
$U(n)$ & the upper bound $\max_q\min(2^{n-q},q-1)$ & Section~\ref{sec:intro}\\
$m$, $q$ & the numbers of states of $X$ and of $\hX$ & Section~\ref{sec:intro}\\
$H(x,j)$, $F_{ij}$ & $\E_xT_j$ and $\E_iT_j-\E_\pi T_j$ & Section~\ref{sec:pot}\\
$\omega$, $\one$, $\delta_x$ & the vector $(w_x)$, the constant $1$, the indicator of $x$ & Section~\ref{sec:pot}\\
$L$, $\mathcal E$ & the Laplacian (loops removed) and the Dirichlet form & Section~\ref{sec:pot}\\
$R(y\lra Z)$ & the effective resistance between $y$ and $Z$ & Section~\ref{sec:pot}\\
$o$, $\Gamma$ & the grounded state and the grounded inverse of $L$ & Section~\ref{sec:pot}\\
$\zeta$, $\bar\zeta$ & $\Gamma\omega$, so $\zeta_i=H(i,o)$, and its $\pi$-mean & Section~\ref{sec:pot}\\
$\kappa_i$, $\varphi_i$ & $\max_jF_{ij}$ and the potential $\kappa_i-F_{i\cdot}$ & Lemma~\ref{lem:pot}\\
$V_{xy}$ & the states reached from $x$ without visiting $y$ & Lemma~\ref{lem:bridge}\\
$S_e$, $w_e$ & the weight beyond an edge $e$ of a tree, and the weight of $e$ & Lemma~\ref{lem:tree}\\
$u$, $v$, $J$, $\hF$ & the lumped pair, the other states, the matrix $F$ of $\hX$ & Section~\ref{sec:merge}\\
$h$, $R$, $\bar h$ & $\Prob_x(T_u<T_v)$, $R(u\lra v)$, the $\pi$-mean of $h$ & Section~\ref{sec:merge}\\
$\beta_i$ & $\varphi_i(u)-\varphi_i(v)=wR(h_i-\bar h)$ & Section~\ref{sec:merge}\\
$\eta$, $\bar\eta$ & the pointwise minimum of $\varphi_i/\beta_i$ and $\varphi_j/\beta_j$, and its $\pi$-mean & Section~\ref{sec:merge}\\
$\pi_*$ & $\min_x\pi_x$ & Remark~\ref{rem:pimin}\\
$N$, $t$, $r$, $\alpha$, $\ell$ & the parameters of the comb & Section~\ref{sec:comb}\\
$a$, $b$, $c$, $z_k$, $y_k$ & the states of the comb: two ends, the heavy leaf, spine, teeth & Section~\ref{sec:comb}\\
$\hat z_k$, $\hat c$ & the lumped states $f(z_k)=f(y_k)$ and $f(z_N)=f(c)$ & Section~\ref{sec:comb}\\
$\sigma_k$, $B$, $S_k$, $\hS_k$ & the constants defined before Proposition~\ref{prop:comb} & Section~\ref{sec:comb}\\
$M$, $\hM$ & $\E_aT_c$ in $X$ and $\hE_a\widehat T_{\hat c}$ in $\hX$ & Section~\ref{sec:comb}\\
$G_Z(x,y)$ & the expected number of visits to $y$ before $T_Z$ from $x$ & Section~\ref{sec:hit}\\
\bottomrule
\end{tabular}
\end{center}


\begin{thebibliography}{9}
\bibitem{Aldous82} D.~J.~Aldous, Some inequalities for reversible Markov chains, \emph{J. London Math. Soc.} (2)
\textbf{25} (1982), 564--576. \href{https://doi.org/10.1112/jlms/s2-25.3.564}{doi:10.1112/jlms/s2-25.3.564}.
\bibitem{AF} D.~Aldous and J.~A.~Fill, \emph{Reversible Markov Chains and Random Walks on Graphs}, unfinished
monograph, 2002, recompiled 2014, \url{https://www.stat.berkeley.edu/~aldous/RWG/book.html}.
\bibitem{HP} J.~Hermon and Y.~Peres, On sensitivity of mixing times and cutoff, \emph{Electron. J. Probab.}
\textbf{23} (2018), paper no.~25, 1--34. \href{https://doi.org/10.1214/18-EJP154}{doi:10.1214/18-EJP154};
arXiv:1610.04357.
\bibitem{LW95} L.~Lov\'asz and P.~Winkler, Efficient stopping rules for Markov chains, in \emph{Proceedings of the
Twenty-Seventh Annual ACM Symposium on Theory of Computing (STOC '95)}, ACM, New York, 1995, 76--82.
\href{https://doi.org/10.1145/225058.225086}{doi:10.1145/225058.225086}.
\bibitem{LW98} L.~Lov\'asz and P.~Winkler, Mixing times, in \emph{Microsurveys in Discrete Probability}, DIMACS Ser.
Discrete Math. Theoret. Comput. Sci. \textbf{41}, Amer. Math. Soc., Providence, RI, 1998, 85--133.
\href{https://doi.org/10.1090/dimacs/041/06}{doi:10.1090/dimacs/041/06}.
\bibitem{LP} R.~Lyons and Y.~Peres, \emph{Probability on Trees and Networks}, Cambridge Series in Statistical and
Probabilistic Mathematics \textbf{42}, Cambridge University Press, New York, 2016.
\href{https://doi.org/10.1017/9781316672815}{doi:10.1017/9781316672815}.

\bibitem{MathBrainRepo} Y.~Shmalo, \emph{Math Brain: public snapshot, code and tests}, 2026,
\url{https://github.com/yspennstate/math-brain-public}.
\bibitem{MathBrainPaper} Y.~Shmalo, \emph{Math Brain: shared, auditable memory for multi-agent mathematical research},
October 2026, \url{https://github.com/yspennstate/math-brain-public/tree/main/papers/math-brain-2026-10-24}.
\bibitem{BrainFramework} Y.~Shmalo, \emph{Brain Framework}, public software repository, 2026,
\url{https://github.com/yspennstate/brain-framework-public}.

\end{thebibliography}
\end{document}
